% !TEX root = ../main.tex
\section{Q02: Transfer Pricing -- under costless og asymmetrisk informasjon}\label{q:02}

\begin{quote}
\emph{Redegør -- med udgangspunkt i forudsætningerne om: a) costless information b) asymmetrisk information -- for begrebet Transfer Pricing (interne afregningspriser). På baggrund heraf ønskes der redegjort for de i pensum nævnte metoder til fastlæggelse af Transfer Pricing, ligesom fordele og ulemper ved metoderne bedes belyst.}
\end{quote}

\subsection*{Svar (kort)}
Transfer pricing er prisen som settes når en divisjon leverer internt til en annen divisjon, og den teoretisk korrekte prisen er lik selgerens \emph{opportunity cost}. Under \emph{costless information} (a) kan toppledelsen sette denne prisen direkte og oppnå konsernoptimal mengde. Under \emph{asymmetrisk informasjon} (b) sitter divisjonene på spesifik viden og har insentiv til å rapportere kostnader strategisk -- ingen prisingsmetode er da førstebeste, og valget mellom market-based, marginal cost, full cost, full cost plus markup og negotiated transfer pricing er en \emph{trade-off mellom decision management (riktig beslutning) og decision control (riktig insentiv)}. [SLIDE]

\subsection*{Relevante teorier (kort)}

\paragraph{\thref{transfer-pricing}{Transfer Pricing}\quad\SOne/\STwo/\SThree}
Oppstår nødvendigvis i konserner med profit-/investmentcentre og interne leveranser. Optimal pris $=$ opportunity cost: marginalkostnad ved ledig kapasitet, eller MC pluss fortrengt dekningsbidrag ved kapasitetsbinding. Påvirker samtidig divisjonens profit (S2), lederens insentiv (S3) og hvilke produksjons- og investeringsbeslutninger som faktisk treffes (S1). De fem metodene i pensum er ulike praktiske tilnærminger til denne ideelle prisen. [SLIDE]

\paragraph{\thref{responsibility-centers}{Responsibility centers}\quad\SOne/\STwo}
Transfer pricing er en direkte konsekvens av valget om å organisere konsernet i \emph{profit centers} eller \emph{investment centers}: når divisjonen har eget P\&L og samtidig leverer internt, må den interne transaksjonen prises. Cost- og expense-centre slipper problemstillingen, fordi de ikke har inntektsside. Transfer pricing er derfor det viktigste praktiske ulempeproblemet ved profitcenter-strukturen. [SLIDE]

\paragraph{\thref{specific-vs-general-knowledge}{Spesifik vs.\ generell viden} \& \thref{principal-agent}{Principal-agent}\quad\SOne/\SThree}
Asymmetrisk informasjon (b) bygger på Jensen \& Mecklings poeng om at lokal kunnskap er kostbar å overføre oppover. Divisjonslederen har spesifik viden om sin egen MC, kapasitet og marked; toppledelsen kan ikke verifisere det. Dette er et prinsipal-agent-problem på tvers av divisjonsgrenser, og er hele grunnen til at transfer pricing er vanskelig i praksis -- ellers kunne toppledelsen bare diktert riktig pris. [SLIDE/BOK]

\paragraph{\thref{decision-management-control}{Decision management vs.\ decision control}\quad\STwo/\SThree}
Samme transfer pris brukes både til \emph{beslutning} (skal vi handle internt eller eksternt? hvor mye?) og til \emph{kontroll} (hvor stor profit får divisjonslederen kreditert, og dermed bonus?). Det som er optimalt for beslutningen (presis MC) er ikke alltid optimalt for kontrollen (gir lederen insentiv til å overrapportere MC). Dette trade-off-et er kjernen i hvorfor ingen enkelt metode dominerer. [SLIDE]

\subsection*{Illustrasjon}

\begin{center}
\renewcommand{\arraystretch}{1.18}
\begin{tabular}{@{}p{3.0cm}p{4.6cm}p{4.8cm}p{2.2cm}@{}}
\toprule
\textbf{Metode} & \textbf{Fordele} & \textbf{Ulemper} & \textbf{Best n\aa r\dots} \\
\midrule
Market-based & $=$ opportunity cost; rettferdig; lite gaming & Krever likvid eksternt marked; pris-volatilitet & Likvid eksternt marked \\
Marginal cost (MC) & Konsernoptimal mengde ved ledig kapasitet & Dekker ikke faste kost; selger taper; MC overrapporteres & Ledig kapasitet \\
Full cost & Enkel; dekker faste kost & $>$ opportunity cost; for lav intern mengde; double markup & Stabil produksjon \\
Full cost + markup & Selger f\aa r margin; bedre selger-insentiv & Forsterker double markup; sterkere gaming-insentiv & Selger trenger insentiv \\
Negotiated & N\ae rmer seg opp.\ cost; bruker lokal info & Tidkrevende; avhengig av forhandlingsstyrke; konflikter & Likeverdige parter \\
\bottomrule
\end{tabular}
\end{center}
\textit{Tabell: Pensums fem metoder for transfer pricing -- fordele, ulemper og forutsetninger. [SLIDE -- BSZ kap.\ 17]}

% Slide-figur fra BSZ kap. 17 (slide 25 ``Decentralized firm with two profit centers and transfer pricing'') finnes som EMF og må konverteres til PDF/PNG for \includegraphics. Når konvertert, aktiveres:
% \begin{figure}[h]
%   \centering
%   \includegraphics[width=0.7\linewidth]{BSZ_kap_17/by_slide/slide_NNN/<filnavn>.pdf}
%   \caption{Decentralized firm with two profit centers and transfer pricing. [SLIDE]}
% \end{figure}

\subsection*{(a) Costless information -- førstebeste-løsning}
\begin{itemize}\setlength{\itemsep}{2pt}
\item Toppledelsen kjenner alle MC-, kapasitets- og etterspørselsfunksjoner.
\item Optimal transfer pris $=$ selgerens \textbf{marginal opportunity cost}: $MC$ ved ledig kapasitet; $MC + $ fortrengt dekningsbidrag ved kapasitetsbinding.
\item Resultat: kjøperdivisjonen utvider intern handel akkurat til $MR_{\text{kj\o per}} = MC_{\text{selger}}$ -- ekvivalent med konsernets MR$=$MC.
\item I dette regimet er valget av \emph{metode} mindre viktig: alle metoder kan kalibreres til samme pris. [SLIDE]
\end{itemize}

\subsection*{(b) Asymmetrisk informasjon -- andrebeste-løsninger}
\begin{itemize}\setlength{\itemsep}{2pt}
\item Divisjonslederne sitter med spesifik viden om kostnad, kapasitet og marked; toppledelsen kan ikke verifisere uten kostnad.
\item \emph{Selger} har insentiv til å \textbf{overrapportere} MC for å presse opp prisen $\Rightarrow$ for lav intern mengde.
\item \emph{Kj\o per} har insentiv til å \textbf{underdrive} betalingsvilje eller bl\aa se opp eksterne alternativer for \aa\ presse prisen ned.
\item Trade-off: presise tall (decision management) vs.\ insentivkompatible tall (decision control) -- jf.\ slide 15 ``samme tall kan ikke alltid brukes til begge formål''.
\item Konsekvens: \emph{ingen} av de fem metodene er førstebeste; valget av metode må veie konsernets viktigste sårbarhet (gaming, double markup, manglende marked, kapasitetsbinding). [SLIDE/BOK]
\end{itemize}

\subsection*{Forutsetninger for hver metode}
\begin{itemize}\setlength{\itemsep}{2pt}
\item \textbf{Market-based:} Eksternt likvid og konkurransedyktig marked finnes; markedspris $\approx$ opportunity cost.
\item \textbf{Marginal cost:} MC kan observeres/verifiseres; ledig kapasitet i selgerdivisjonen.
\item \textbf{Full cost:} Stabil produksjonsmiks; faste kostnader kan allokeres meningsfullt; selger trenger dekning.
\item \textbf{Full cost + markup:} Som full cost, men selger må sikres positiv margin (motivasjon).
\item \textbf{Negotiated:} Likeverdig forhandlingsstyrke; toppledelsen er villig til å eskalere når forhandlingene strander; transaksjonsfrekvens lav nok til at forhandlingstid er forsvarlig.
\end{itemize}

\subsection*{Real-world eksempler}

\textbf{Eks.~1 -- Equinor (NO).}\quad [EKS]
Letingsdivisjonen leverer olje og gass internt til Marketing \& Trading-divisjonen, som videreselger på spotmarkedet. Equinor bruker \emph{market-based} transfer pricing knyttet til Brent / NBP-priser. Dette fungerer fordi det finnes et likvid og konkurransedyktig eksternt marked: markedspris $\approx$ opportunity cost. Resultat: minimal intern friksjon, og leting evalueres mot reell verdiskaping i markedet -- ikke mot interne forhandlingsutfall.

\textbf{Eks.~2 -- Mærsk (DK), sammenligning med Eks.~1.}\quad [EKS]
Mærsks Ocean-divisjon leier konsernspesifikke containere og terminaltjenester av Logistics \& Services. Det finnes ikke noe likvid eksternmarked for nøyaktig samme tjeneste, så market-based er ikke et alternativ. Mærsk har historisk brukt en blanding av \emph{cost-plus} og \emph{negotiated} transfer pricing. Sammenlignet med Equinor: samme prinsipielle problem (intern leveranse mellom profitcentre), men ulik metode fordi forutsetningen om et eksternt marked brytes. Mærsk har dokumenterte interne konflikter rundt prising av disse tjenestene -- klassisk symptom på asymmetrisk informasjon (b).

\subsection*{Kobling til Mærsk Power-to-X}
Mærsks Power-to-X-satsing skal produsere grønn metanol og grønn hydrogen som internt drivstoff til shippingdivisjonens nye metanol-skip. Dette er en lærebokssituasjon for Q02: [CASE]
\begin{itemize}\setlength{\itemsep}{2pt}
\item \textbf{Eksternt marked mangler:} Grønn metanol har enn\aa\ ikke et likvid spotmarked $\Rightarrow$ market-based faller bort.
\item \textbf{MC-pricing problematisk:} PtX-anlegget har enorme faste kapital- og strømkontraktskostnader; ren MC dekker ikke disse, og PtX-divisjonen vil tape penger på hver intern leveranse.
\item \textbf{Full cost / cost-plus:} Mest sannsynlige utfall, men risikerer at shippingdivisjonen heller velger fossilt drivstoff fra eksterne leverandører fordi intern grønn metanol blir for dyr -- en intern \emph{double markup}-effekt som motvirker hele Power-to-X-strategien.
\item \textbf{Negotiated:} Sannsynlig overgangsløsning; eskaleres til konsernledelsen som strategisk pris under markedet for å sikre etterspørsel i oppstartsfasen.
\item \textbf{Strategisk subsidie:} Mærsk kan måtte sette transfer prisen \emph{under} full cost (dual pricing eller bevisst tap) for å oppnå konsernets klimastrategiske mål -- viser hvordan transfer pricing også er et strategisk virkemiddel, ikke bare et regnskapsteknisk valg.
\end{itemize}

\subsection*{Fallgruber}
\begin{itemize}\setlength{\itemsep}{2pt}
\item Glemme \emph{opportunity cost}-prinsippet -- nevn alltid at dette er det teoretisk korrekte, og at de fem metodene er praktiske approksimasjoner.
\item Ikke skille tydelig mellom (a) costless og (b) asymmetrisk informasjon -- spørsmålet ber eksplisitt om begge.
\item Glemme \emph{double markup} ved full cost / cost-plus.
\item Glemme \emph{kapasitetsbetingelsen} for MC-pricing (MC er bare opp.\ cost ved ledig kapasitet).
\item Behandle metodene isolert -- alltid si \emph{når} hver metode er best, ikke bare fordele/ulemper.
\item Hoppe over koblingen til \emph{decision management vs.\ decision control}-trade-off-et -- sensorene vil h\o re den.
\item Bare ett eksempel; spørsmålet inviterer til sammenligning av ulike metoder i ulike kontekster.
\end{itemize}

\subsection*{Håndskrivbar disposisjon (15-min forberedelse)}
\begin{itemize}\setlength{\itemsep}{2pt}
\item \textbf{Tese:} Transfer pris $=$ opportunity cost. (a) costless $\Rightarrow$ f\o rstebeste; (b) asymmetrisk $\Rightarrow$ trade-off, ingen metode dominerer.
\item \thref{transfer-pricing}{Transfer pricing}: oppst\aa r i profit-/investmentcentre med intern handel; p\aa virker S1+S2+S3 samtidig.
\item \thref{responsibility-centers}{Centertyper}: TP er ulempen ved profitcenter-strukturen.
\item (a) \textbf{Costless info}: TP $=$ MC (ledig kap.) eller MC $+$ fortrengt DB (kap.\ binding) $\Rightarrow$ MR$=$MC for konsernet.
\item (b) \textbf{Asymmetrisk}: \thref{specific-vs-general-knowledge}{spesifik viden} $+$ \thref{principal-agent}{prinsipal-agent} $\Rightarrow$ selger overrapporterer MC, kj\o per underdriver WTP.
\item \thref{decision-management-control}{Decision management vs.\ control}: presise tall $\neq$ insentivkompatible tall.
\item \textbf{5 metoder} (tabell): market-based / MC / full cost / cost-plus / negotiated. Kort fordel + ulempe + ``best n\aa r\dots''.
\item \textbf{Double markup}-fellen ved full cost / cost-plus.
\item Eks.\ 1: Equinor (Brent-pris, market-based, fungerer). Eks.\ 2: M\ae rsk Ocean / L\&S (cost-plus + negotiated, friksjon).
\item \textbf{M\ae rsk PtX}: ingen marked $\Rightarrow$ cost-plus / negotiated / strategisk subsidie -- klassisk asymmetrisk-info-case.
\item S\o yle: \STwo\ prim\ae rt; ogs\aa\ \SOne\ og \SThree.
\end{itemize}
